% HojaTituloPAE contiene toda la informaci\'{o} acerca de la presentaci\'{o}n del Trabajo de Grado

\begin{center}
\begin{figure}
\centering
\epsfig{file=HojaTitulo/Fig_HojaTitulo/LogoUV.eps,scale=0.4}
\end{figure}

\thispagestyle{empty}
\vspace*{2.0cm}

%Reemplace este t\'{\i}tulo por el de su Trabajo de Grado
\textbf{\huge Dise\~{n}o y validaci\'{o}n de muestreo de c\'{i}tricos para detecci\'{o}n de enfermedades en viveros del Valle del Cauca}\\[5.0cm]

%Aqui va los nombres y apellidos completos de los autores
\Large\textbf{Kevin Steven Garc\'{i}a Chica}\\ 
\Large\textbf{Jose Alejandro Vargas Franco}\\[4.0cm]
\small Universidad del Valle\\
Facultad de Ingenier\'{\i}a, Escuela de Estad\'{\i}stica\\
Santiago de Cali, Colombia\\
%Modifique fecha de ser necesario
2019\\
\end{center}

\newpage{\pagestyle{empty}\cleardoublepage}

\newpage

\begin{center}
\thispagestyle{empty}
\vspace*{0cm}

%Reemplace este t\'{\i}tulo por el de su Trabajo de Grado
\textbf{\huge Dise\~{n}o y validaci\'{o}n de muestreo de c\'{i}tricos para detecci\'{o}n de enfermedades en viveros del Valle del Cauca}\\[3cm]

%Aqui va los nombres y apellidos completos de los autores
\Large\textbf{Kevin Steven Garc\'{i}a Chica}\\ 
\Large\textbf{Jose Alejandro Vargas Franco}\\[3.0cm]
\small Trabajo de grado presentado como requisito parcial para optar al t\'{\i}tulo de:\\
%Modifique seg\'{u}n el g\'{e}nero
\textbf{Estad\'{\i}stico}\\[2.5cm]
%Modifique seg\'{u}n el g\'{e}nero
Director(a):\\
Jose Rafael Tovar\\
%Modifique seg\'{u}n el g\'{e}nero
Codirector(a):\\
T\'{\i}tulo y nombre del codirector(a)\\[3.5cm]
Universidad del Valle\\
Facultad de Ingenier\'{\i}a, Escuela de Estad\'{\i}stica\\
Santiago de Cali, Colombia\\
%Modifique fecha de ser necesario
2019\\
\end{center}

\newpage

\thispagestyle{empty} \textbf{}\normalsize
\\\\\\
\textbf{Dedicatoria o un lema}\\[4.0cm]

\begin{flushright}
\begin{minipage}{8cm}
\noindent
\small
Su uso es opcional y cada autor podr\'{a} determinar la distribuci\'{o}n del texto en la p\'{a}gina. Se sugiere esta presentaci\'{o}n en donde el autor dedica su trabajo en forma especial a personas y/o entidades. Por Ejemplo:\\[1.0cm]
A mis padres\\[1.0cm]\\
\'{o}\\[1.0cm]
La preocupaci\'{o}n por el hombre y su destino siempre debe ser el inter\'{e}s primordial de todo esfuerzo t\'{e}cnico. Nunca olvides esto entre tus diagramas y ecuaciones.\\\\
Albert Einstein\\
\end{minipage}
\end{flushright}

\newpage{\pagestyle{empty}\cleardoublepage}

\newpage

\thispagestyle{empty} \textbf{}\normalsize
\\\\\\
\textbf{\LARGE Agradecimientos}\\\\
Esta secci\'{o}n es opcional, en ella el autor agradece a las personas o instituciones que colaboraron en la realizaci\'{o}n del trabajo de grado. Si se incluye esta secci\'{o}n, deben aparecer los nombres completos, los cargos y su aporte al documento.\\

Agradecemos a la Universidad Nacional de Colombia por facilitar el uso de su plantilla institucional para la elaboraci\'{o}n de trabajos de grado, la cual fue tomada como gu\'{\i}a para la construcci\'{o}n de esta.\\

A \'{A}ngela Cadena, estudiante del Programa Acad\'{e}mico de Estad\'{\i}stica, quien con su colaboraci\'{o}n hizo posible la escritura y adaptaci\'{o}n de este documento.

\newpage{\pagestyle{empty}\cleardoublepage}

\newpage

\textbf{\LARGE Resumen}
\addcontentsline{toc}{chapter}{\numberline{}Resumen}\\\\
El resumen es una presentaci\'{o}n abreviada y precisa (la NTC 1486 de 2008 recomienda revisar la norma ISO 214 de 1976). Se debe usar una extensi\'{o}n m\'{a}xima de 12 renglones. Se recomienda que este resumen sea anal\'{\i}tico, es decir, que sea completo, con informaci\'{o}n cuantitativa y cualitativa, generalmente incluyendo los siguientes aspectos: objetivos, dise\~{n}o, lugar y circunstancias (u objetivo del estudio), principales resultados, y conclusiones. Al final del resumen se deben usar palabras claves tomadas del texto (m\'{\i}nimo 3 y m\'{a}ximo 7 palabras), las cuales permiten la recuperaci\'{o}n de la informaci\'{o}n.\\[0.5cm]
\textbf{\small Palabras clave: palabras clave en espa\~{n}ol (m\'{a}ximo 10 palabras, preferiblemente seleccionadas de las listas internacionales que permitan el indizado cruzado)}\\[1cm]

\textbf{\LARGE Abstract}\\\\
Es el mismo resumen pero traducido al ingl\'{e}s. Se debe usar una extensi\'{o}n m\'{a}xima de 12 renglones. Al final del Abstract se deben traducir las anteriores palabras clave tomadas del texto (m\'{\i}nimo 3 y m\'{a}ximo 7 palabras), llamadas keywords. Es posible incluir el resumen en otro idioma diferente al espa\~{n}ol o al ingl\'{e}s, si se considera como importante dentro del tema tratado en la investigaci\'{o}n, por ejemplo: un trabajo dedicado a problemas ling\"{u}\'{\i}sticos del mandar\'{\i}n seguramente estar\'{\i}a mejor con un resumen en mandar\'{\i}n.\\[0.5cm]
\textbf{\small Keywords: palabras clave en ingl\'{e}s (m\'{a}ximo 10 palabras, preferiblemente seleccionadas de las listas internacionales que permitan el indizado cruzado)}\\

\newpage{\pagestyle{empty}\cleardoublepage}